\slajdid{02-s023b}
\begin{frame}[label=02-s023b]{Poređenje resursa dve direktne realizacije}
% element: 02-s023b:e01
\Rezultat{Ista funkcija, zapisana kao zbir proizvoda ili proizvod zbirova, po pravilu zahteva \alert{različit broj resursa}.}
\par\vspace{4mm}
% element: 02-s023b:e02
\begin{columns}[T,onlytextwidth]
\begin{column}{.47\textwidth}
\setlength{\abovedisplayskip}{4pt}\setlength{\abovedisplayshortskip}{4pt}
\Oznaka{Zbir proizvoda}\vspace{2mm}
Tri troulazna I kola i jedno troulazno ILI kolo:
\[3\cdot3+1\cdot3=\textcolor{DEcrvena}{12\text{ ulaza}}.\]
\end{column}
\begin{column}{.47\textwidth}
\setlength{\abovedisplayskip}{4pt}\setlength{\abovedisplayshortskip}{4pt}
\Oznaka{Proizvod zbirova}\vspace{2mm}
Pet troulaznih ILI kola i jedno petoulazno I kolo:
\[5\cdot3+1\cdot5=\textcolor{DEcrvena}{20\text{ ulaza}}.\]
\end{column}
\end{columns}
\par\vspace{3mm}
% element: 02-s023b:e03
U ovom primeru \alert{zbir proizvoda zahteva manje resursa}.
\par\vspace{2mm}
Koji oblik će biti jednostavniji zavisi od \alert{broja i rasporeda nula i jedinica} u koloni funkcije \(F\) u funkcionalnoj tabeli.
\end{frame}
